% =========================================================
% PAGE 27: TABLES INDEX (Tables A1 & A2) - ARABIC
% =========================================================

% 1. Geometry
\GeomStandard


% ---------------------------------------------------------
% PAGE TITLE
% ---------------------------------------------------------
\noindent
{\FontPageTitle\textbf{\color{ecestitle} \ECESContent{title_tables_index_ar}{ملحق الجداول}}}
\vspace{0.5em}
\label{sec:ar_appendix}
% ---------------------------------------------------------
% TABLE A1 (t1.png)
% ---------------------------------------------------------
\begin{figure}[h!]
    \centering
    % Constrain height to 38% of text height to ensure two tables fit
    \includegraphics[width=\linewidth, height=0.38\textheight, keepaspectratio]{t1.png}
\end{figure}

\vspace{0.2em}

% ---------------------------------------------------------
% TABLE A2 (t2.png)
% ---------------------------------------------------------
\begin{figure}[h!]
    \centering
    % Constrain height to 38% of text height
    \includegraphics[width=\linewidth, height=0.38\textheight, keepaspectratio]{t2.png}
\end{figure}

% ---------------------------------------------------------
% FOOTNOTES
% ---------------------------------------------------------
\vfill
\noindent
{\FontNote
\ECESContent{tables_footnotes_pg27_ar}{1 تمثل الأرقام نسبة مئوية من إجمالي الردود. ارتفاع وثبات وانخفاض قد لا يصل مجموعهم إلى 100 بسبب التقريب.\\
2 يساوي المتوسط البسيط لمؤشرات المتغيرات. طريقة حساب المؤشر تم توضيحها في الملحق.}
}

\clearpage

% =========================================================
% PAGE 28: TABLES APPENDIX (Tables A3 & A4) - ARABIC
% =========================================================

% 1. Geometry (Persist from previous page)



% ---------------------------------------------------------
% PAGE TITLE
% ---------------------------------------------------------
\noindent
{\FontPageTitle\textbf{\color{ecestitle} \ECESContent{title_tables_appendix_ar}{ملحق الجداول}}}

\vspace{0.5em}

% ---------------------------------------------------------
% SUBTITLE
% ---------------------------------------------------------
\noindent
{\FontTableTitle
\textbf{\ECESContent{title_table3_ar}{الجدول رقم 3: نتائج الاستبيان، مؤشر الأداء السابق لكافة الشركات، وفقا لحجمها (ابريل - مايو - يونيو 2025)}}
}

\vspace{0.5em}

% ---------------------------------------------------------
% TABLE A3 (t3.png)
% ---------------------------------------------------------
\begin{figure}[h!]
    \centering
    % Table A3 - Standard width
    \includegraphics[width=\linewidth, height=0.35\textheight, keepaspectratio]{t3.png}
\end{figure}

\vspace{0.8em}

% ---------------------------------------------------------
% SUBTITLE FOR TABLE A4
% ---------------------------------------------------------
\noindent
{\FontTableTitle
\textbf{\ECESContent{title_table4_ar}{الجدول رقم 4: نتائج الاستبيان، مؤشر التوقعات لكافة شركات العينة، وفقا لحجمها (ابريل - أغسطس - سبتمبر 2025)}}
}

\vspace{0.5em}

% ---------------------------------------------------------
% TABLE A4 (t4.png)
% ---------------------------------------------------------
\begin{figure}[h!]
    \centering
    % Table A4 - Standard width
    \includegraphics[width=\linewidth, height=0.35\textheight, keepaspectratio]{t4.png}
\end{figure}

% ---------------------------------------------------------
% FOOTNOTES
% ---------------------------------------------------------
\vfill
\noindent
{\FontNote
\ECESContent{tables_footnotes_pg28_ar}{1 تمثل الأرقام نسبة مئوية من إجمالي الردود. ارتفاع وثبات وانخفاض قد لا يصل مجموعهم إلى 100 بسبب التقريب.\\
2 يساوي المتوسط البسيط لمؤشرات المتغيرات. طريقة حساب المؤشر تم توضيحها في الملحق.}
}

\restoregeometry
\clearpage